% Replacement for Table 2 (tab:kan_config) in Section 3.5 of the paper.
% The forecasting column is the original table unchanged; the third column
% (red) carries the approximation-study settings so that Section 4.2 needs
% no table of its own. The "Grid range / update" row is new (red): both
% values are PyKAN defaults in the forecasting code (driver.py, KAN(...) and
% net.fit(dataset, opt="LBFGS", steps=10)) and explicit arguments in
% approx_study.py (grid_range=[0,1], update_grid=False).
% Section 3.5's \subsection{KAN Configuration} was given \label{sec:kan_config_fc}.
\begin{table}[h]
\centering
\caption{KAN implementation specification for the forecasting study (Section~\ref{sec:kan_config_fc}) and, {\JOY in red}, the approximation study (Section~\ref{approx_kan}); \emph{same} means the forecasting value is used}
\label{tab:kan_config}
\footnotesize
\begin{tabular}{lll}
\hline
\textbf{Parameter} & \textbf{Forecasting} & {\JOY\textbf{Approximation}} \\
\hline
Framework & PyTorch + PyKAN library & {\JOY same} \\
Architecture & \texttt{KAN(width=[input, \(N_h\), \(H\)])} & {\JOY \texttt{[1,1]} and \texttt{[1,\(h\),1]}} \\
Input representation & Flattened: \(L \cdot |\mathcal{F}|\) & {\JOY scalar \(t\in[0,1]\)} \\
Hidden neurons (\(N_h\)) & \(3 \times \lfloor |\text{train set}| / 50 \rfloor\) & {\JOY none; \(h\in\{1,1,2,5,10\}\) by budget} \\
Grid size & 3 & {\JOY \([1,h,1]\): 3; \([1,1]\): \(G\in\{6,21,51,111,231\}\)} \\
{\JOY Grid range / update} & {\JOY \([-1,1]\) (PyKAN default), adaptive} & {\JOY \([0,1]\), fixed (no update)} \\
\(k\) (B-spline degree) & 3 & {\JOY same} \\
Optimiser & L-BFGS (full-batch, second-order) & {\JOY same} \\
Training steps & 10 L-BFGS iterations & {\JOY 80} \\
Regularisation & PyKAN defaults only & {\JOY same} \\
Random seed & 0 (fixed) & {\JOY \(\{0,1,2\}\)} \\
Stopping criteria & Fixed 10 steps (no early stopping) & {\JOY fixed 80 steps (no early stopping)} \\
\hline
\end{tabular}
\end{table}
